\documentclass{situationist-preprint}

\title{The Self-Undermining Objection Answered}
\msid{EPI-07-b}
\msfamily{EPI-07}
\msdiscipline{Epistemology}
\mstype{Companion manuscript}
\mskeywords{peer disagreement; conciliationism; self-undermining; higher-order evidence}

\begin{document}
\maketitle

\begin{abstract}
Conciliatory views of disagreement face the objection that, since many of their proponents' peers reject them, they recommend their own rejection. We show that the moderate conciliationism of EPI-07-a, applied to disagreement about itself, recommends some loss of confidence in itself but not its abandonment: the conciliationist ends strictly between her credence and her opponent's, closer to her own. We then answer two ways of pressing the objection further. Repeated encounters with the same opponents do not drive her credence down, because the rational credence is fixed by the evidence and repeated disclosure adds none; and a view that recommends lowered confidence in itself is not thereby incoherent.
\end{abstract}

\section{The objection}

\citet{Elga2010} argued that conciliatory views are self-undermining: there is peer disagreement about whether to conciliate, so a conciliationist who meets a steadfast peer should, by her own view, lose confidence in conciliationism, and in the limit be led to reject it. A view about how to respond to evidence that recommends its own rejection in some circumstances is, Elga argued, incoherent, and the only escape is to make conciliationism dogmatic about itself.

EPI-07-a holds that the rational credence after disclosure is $r = \lambda f + (1 - \lambda)\frac{a+b}{2}$, where $f$ is the credence the first-order evidence supports and $0 < \lambda < 1$ (its Principles~P1--P3), and that each party can follow the moderate policy of moving from her credence $x$ to $\frac{1+\lambda}{2}x + \frac{1-\lambda}{2}y$ when her peer's is $y$ (Definition~3.1).

\section{Self-application}

Let $C$ be the proposition that moderate conciliationism is correct. Let $A$ have credence $a$ in $C$ and her peer $B$ credence $b < a$.

\begin{proposition}\label{p:loss}
Applied to the disagreement about $C$, the moderate policy moves $A$ to a credence $o_A$ with $b < o_A < a$ and $a - o_A < o_A - b$.
\end{proposition}

\begin{proof}
This is Theorem~3.2 of EPI-07-a with the proposition in dispute taken to be $C$.
\end{proof}

So the view recommends a loss of confidence in itself, which is the correct verdict: the disagreement of a peer is evidence that one may have misjudged the arguments, and that is as true of arguments about disagreement as of any others. It does not recommend adopting the opponent's credence, nor giving the opponent's view more credence than one's own, and it is in this sense that it does not recommend its own abandonment.

The objection may be pressed by iteration: if $A$ meets $B$ again, does she conciliate again, and so converge on $B$'s credence?

\begin{proposition}\label{p:fixed}
Repeated disclosure of the same disagreement leaves the rational credence unchanged.
\end{proposition}

\begin{proof}
By Principle~P3 of EPI-07-a, the rational credence is a function of $f$, $a$ and $b$, where $a$ and $b$ are the parties' credences on the first-order evidence before disclosure. A second disclosure of the same credences changes none of these, so $r$ is unchanged.
\end{proof}

Iterating the policy would count the same higher-order evidence twice. The policy is a rule for responding to a disclosure, not a dynamics to be run until convergence. The same holds when the steadfast peers are many: the higher-order component then combines the credences of all parties, once each, and the first-order component retains positive weight; however numerous the opponents, the view never prescribes taking over their credence, since $\lambda f$ remains in the sum.

\section{Coherence}

\begin{objection}
A rule that recommends less than full confidence in itself gives inconsistent advice: it tells the agent to follow it, and also, by lowering her confidence in it, to give weight to rules that conflict with it \citep{Elga2010}.
\end{objection}

\begin{reply}
The rule says what credence is rational given the evidence, and so it applies to credence in itself as to anything else. Its output for $C$ is a credence below one; it does not follow that the agent should follow some other rule. An agent who is fairly confident that a rule is correct, and has no better candidate, follows it while acknowledging that she may be wrong, as she does with every substantive principle she accepts. The demand that a rule of evidence be certain of itself is a demand that no substantive epistemic principle meets, and \citet{Christensen2013} has argued that the self-undermining charge, so understood, applies to any principle of epistemic modesty. What would be incoherent is a rule whose recommendation for its own credence led to recommendations for first-order credences that it itself forbids. By Proposition~\ref{p:fixed} the moderate view yields a single determinate credence for every proposition, $C$ included, and does not.
\end{reply}

\section{Conclusion}

Moderate conciliationism, applied to disagreement about itself, recommends a moderate loss of confidence in itself. Because it weights first-order evidence and counts each disclosure once, the loss does not compound into abandonment, and the view's self-application is consistent.

\begin{thebibliography}{9}\small
\bibitem[Christensen(2013)]{Christensen2013} Christensen, D. (2013). Epistemic modesty defended. In D.~Christensen and J.~Lackey (Eds.), \emph{The Epistemology of Disagreement: New Essays}. Oxford University Press.
\bibitem[Elga(2010)]{Elga2010} Elga, A. (2010). How to disagree about how to disagree. In R.~Feldman and T.~A. Warfield (Eds.), \emph{Disagreement} (pp.~175--186). Oxford University Press.
\end{thebibliography}

\end{document}
